% generated by paper/appendix4_table1.py -- do not edit by hand
\begingroup\footnotesize\setlength{\tabcolsep}{4pt}
\begin{longtable}{@{}r r >{\raggedright\arraybackslash}p{0.44\linewidth} >{\raggedright\arraybackslash}p{0.30\linewidth}@{}}
\caption{\textbf{\AppTabFactors{}. The factors of the factor model under oblimin and varimax rotation.} One row per factor of the factor model with 25 factors, fitted to the answers. A factor's core items are the items with its largest loadings, found by the participation-ratio rule used for the semantic groups (Appendix~\AppGroups{}), and each factor is named by the constructs of all its core items, with the number of core items per construct, largest first. Factors are numbered in the order the rotation returns them, so a number refers to different factors under the two rotations, and the factor numbers in the text are those of the oblimin rotation. Each stand-alone question is counted under the study it was taken from. The questionnaire and source abbreviations are spelled out in Appendix~\AppGroups{}, and the core items of every factor are listed in Supplementary File~\SuppItems{}.}\label{apptab:factors}\\
\toprule Factor & Core items & Constructs of the core items & Questionnaires of the core items \\ \midrule
\endfirsthead \toprule Factor & Core items & Constructs of the core items & Questionnaires of the core items \\ \midrule \endhead
\bottomrule \endfoot
\multicolumn{4}{@{}p{0.97\linewidth}}{\textbf{A. Oblimin rotation (the factors may correlate), the paper's F numbers} \quad 25 factors, 17 on 1 construct, 6 on 2, 2 on 3 or more, 8 with core items from more than 1 questionnaire} \\[2pt]
F1 & 3 & Extraversion 2, Loneliness 1 & BFI-2-XS, ULS-8 \\
F2 & 6 & Daily spiritual experience 5, Religious participation 1 & DSES, NHS \\
F3 & 3 & Health care discrimination 3 & DMS \\
F4 & 3 & Everyday discrimination 3 & EDS \\
F5 & 4 & Neighborhood disorder and safety 4 & NDS \\
F6 & 5 & Childhood adversity 5 & ACE \\
F7 & 5 & Loneliness 5 & ULS-8 \\
F8 & 4 & Perceived social support 4 & mMOS-SS \\
F9 & 6 & Adult ADHD symptoms 4, Conscientiousness 2 & ASRS, BFI-2-XS \\
F10 & 4 & Intimate partner or sexual violence 4 & MHQ \\
F11 & 8 & Psychotic experiences 4, Suicidal ideation and self-harm 4 & MHQ, PHQ-9, SITBI, WMH-CIDI \\
F12 & 5 & Neighborhood social cohesion 4, Neighborhood disorder and safety 1 & NDS, SCNS \\
F13 & 2 & Food insecurity 2 & HVS \\
F14 & 3 & Depressive symptoms 2, Suicidal ideation and self-harm 1 & MHQ, SITBI \\
F15 & 5 & Substance use 5 & PATH, TUS-CPS and MVP \\
F16 & 3 & Childhood adversity 3 & ACE \\
F17 & 2 & Perceived social support 2 & mMOS-SS \\
F18 & 4 & Adult ADHD symptoms 2, Anxiety symptoms 1, Depressive symptoms 1 & ASRS, GAD-7, PHQ-9 \\
F19 & 6 & Perceived stress 6 & PSS-10 \\
F20 & 4 & Perceived stress 4 & PSS-10 \\
F21 & 5 & Agreeableness 3, Negative emotionality 1, Open-mindedness 1 & BFI-2-XS \\
F22 & 4 & Health care discrimination 4 & DMS \\
F23 & 4 & Anxiety symptoms 4 & GAD-7 \\
F24 & 5 & Depressive symptoms 5 & PHQ-9 \\
F25 & 3 & Neighborhood disorder and safety 3 & NDS \\
\addlinespace[4pt]
\multicolumn{4}{@{}p{0.97\linewidth}}{\textbf{B. Varimax rotation (the factors are uncorrelated), a check} \quad 25 factors, 16 on 1 construct, 6 on 2, 3 on 3 or more, 10 with core items from more than 1 questionnaire} \\[2pt]
F1 & 10 & Anxiety symptoms 6, Depressive symptoms 3, Self-rated mental health 1 & GAD-7, PHQ-9, PROMIS \\
F2 & 6 & Daily spiritual experience 5, Religious participation 1 & DSES, NHS \\
F3 & 7 & Perceived social support 7 & mMOS-SS \\
F4 & 4 & Health care discrimination 4 & DMS \\
F5 & 9 & Neighborhood disorder and safety 9 & NDS \\
F6 & 3 & Childhood adversity 3 & ACE \\
F7 & 4 & Everyday discrimination 4 & EDS \\
F8 & 2 & Food insecurity 2 & HVS \\
F9 & 2 & Perceived social support 2 & mMOS-SS \\
F10 & 6 & Adult ADHD symptoms 4, Conscientiousness 2 & ASRS, BFI-2-XS \\
F11 & 5 & Neighborhood social cohesion 4, Neighborhood disorder and safety 1 & NDS, SCNS \\
F12 & 6 & Perceived stress 6 & PSS-10 \\
F13 & 3 & Depressive symptoms 2, Suicidal ideation and self-harm 1 & MHQ, SITBI \\
F14 & 4 & Adult ADHD symptoms 2, Manic experiences 1, Anxiety symptoms 1 & ASRS, GAD-7, MHQ \\
F15 & 5 & Loneliness 5 & ULS-8 \\
F16 & 3 & Health care discrimination 3 & DMS \\
F17 & 4 & Intimate partner or sexual violence 4 & MHQ \\
F18 & 4 & Extraversion 3, Loneliness 1 & BFI-2-XS, ULS-8 \\
F19 & 5 & Neighborhood disorder and safety 5 & NDS \\
F20 & 5 & Agreeableness 3, Negative emotionality 1, Manic experiences 1 & BFI-2-XS, MHQ \\
F21 & 5 & Childhood adversity 5 & ACE \\
F22 & 4 & Depressive symptoms 4 & PHQ-9 \\
F23 & 4 & Perceived stress 4 & PSS-10 \\
F24 & 8 & Psychotic experiences 4, Suicidal ideation and self-harm 4 & MHQ, PHQ-9, SITBI, WMH-CIDI \\
F25 & 5 & Substance use 5 & PATH, TUS-CPS and MVP \\
\addlinespace[4pt]
\end{longtable}
\endgroup
